\subsection{WALLS AND FACES}

DEFINITION 3. Let C be a chamber of E. A face of C is a facet contained in the closure of C whose support is a hyperplane. A wall of C is a hyperplane that is the support of a face of C.

Every wall of C belongs to H. Prop. 4 shows that a hyperplane \(L \in H\) is a wall of C if and only if \(C \neq D_{\mathfrak{H} - \{L\}}(C)\) . Moreover, every wall of C is the support of a single face of C.

PROPOSITION 8. Every hyperplane H belonging to H is the wall of at least one chamber.

By Prop. 7, (iii), there exists a point \(a\) of \(\mathbf{H}\) that does not belong to any hyperplane \(\mathbf{H}' \neq \mathbf{H}\) of \(\mathfrak{H}\) ; by Prop. 6, there exists a chamber \(\mathbf{C}\) such that \(a \in \overline{\mathbf{C}}\) ; Prop. 4 then shows that \(\mathbf{H}\) is a wall of \(\mathbf{C}\) .

PROPOSITION 9. Let C be a chamber and M the set of walls of C. Then \(\mathrm{C} = \mathrm{D}_{\mathfrak{M}}(\mathrm{C})\) and every subset L of H such that \(\mathrm{C} = \mathrm{D}_{\mathfrak{L}}(\mathrm{C})\) contains M. A subset F of \(\overline{C}\) is a facet if and only if it is a facet of E relative to the family M.

\begin{enumerate}
\def\labelenumi{\alph{enumi})}
\item
  Let \(\mathfrak{L}\) be a subset of \(\mathfrak{H}\) such that \(C = D_{\mathfrak{L}}(C)\) . Consider a hyperplane \(L\) belonging to \(\mathfrak{H}\) but not to \(\mathfrak{L}\) ; let \(\mathfrak{N}\) be the set of hyperplanes \(H \neq L\) belonging to \(\mathfrak{H}\) . Then \(\mathfrak{L} \subset \mathfrak{N}\) , hence \(C = D_{\mathfrak{M}}(C)\) , and \(L\) does not meet \(D_{\mathfrak{M}}(C)\) . By the implication (i) \(\Longrightarrow\) (iii) in Prop. 4, the hyperplane \(L\) is not a wall of \(C\) . Consequently, every wall of \(C\) belongs to \(\mathfrak{L}\) .
\item
  We assume that \(C = D_{\mathfrak{L}}(C)\) . Let H be a hyperplane belonging to L that is not a wall of C, and put \(L' = L - \{H\}\) . By the implication (iii) \(\Longrightarrow\) (i) in Prop. 4, the convex set \(D_{\mathfrak{L}'}(C)\) does not meet H, so \(D_{\mathfrak{L}'}(C) \subset D_{H}(C)\) and \(C = D_{\mathfrak{L}'}(C)\) . If F is a finite subset of L that does not contain any wall of C, we conclude by induction on the cardinal of F that \(C = D_{\mathfrak{L} - \mathfrak{F}}(C)\) .
\item
  Let \(a\) be a point of C; clearly, \(\mathrm{C} \subset \mathrm{D}_{\mathfrak{M}}(a)\) . Let \(a'\) be a point of \(\mathrm{D}_{\mathfrak{M}}(a)\) ; since the closed segment \([aa']\) is compact, the set \(\mathfrak{F}\) of hyperplanes \(\mathrm{H} \in \mathfrak{H}\) that meet \([aa']\) is finite. Since \(a\) and \(a'\) are strictly on the same side of every wall of C, no wall of C belongs to \(\mathfrak{F}\) ; by \(b\) , we have \(\mathrm{C} = \mathrm{D}_{\mathfrak{H}-\mathfrak{F}}(\mathrm{C})\) . Since \(a' \in \mathrm{D}_{\mathfrak{H}-\mathfrak{F}}(a)\) , we have \(a' \in \mathrm{C}\) . We have therefore proved that \(\mathrm{D}_{\mathfrak{M}}(a) \subset \mathrm{C}\) , which establishes the first part of the proposition.
\item
  To prove the last assertion of the proposition, it clearly suffices to show that a subset F of \(\overline{C}\) that is a facet of E relative to M is a facet of E relative to H, or that every hyperplane \(H \in H\) that meets F contains F. So let H be a hyperplane that meets F but does not contain it. Since F is open in its affine support, it is not completely on one side of H. It follows that \(\overline{C}\) is not completely on one side of H and hence that the hyperplane H does not belong to H, which completes the proof.
\end{enumerate}

Remarks. 1) It follows from formula (6) and Prop. 9 that the closure of a chamber C is the intersection of the closed half-spaces that are bounded by a wall of C and contain C.

\begin{enumerate}
\def\labelenumi{\arabic{enumi})}
\setcounter{enumi}{1}
\tightlist
\item
  Let \(\mathbf{F}\) be a facet whose support is a hyperplane \(\mathbf{L}\) ; we shall show that \(\mathbf{F}\) is a face of two chambers. Let \(\mathfrak{N}\) be the set of hyperplanes \(\mathrm{H} \neq \mathrm{L}\) belonging to \(\mathfrak{H}\) ; put \(\mathrm{A} = \mathrm{D}_{\mathfrak{N}}(\mathrm{F})\) and denote by \(\mathrm{D}^{+}\) and \(\mathrm{D}^{-}\) the two open half-spaces bounded by \(\mathrm{L}\) . The set \(\mathbf{A}\) is open and contains \(\mathbf{F} \subset \mathbf{L}\) , and since every point of \(\mathbf{L}\) is in the closure of \(\mathrm{D}^{+}\) and \(\mathrm{D}^{-}\) , the sets \(\mathrm{C}^{+} = \mathrm{A} \cap \mathrm{D}^{+}\) and \(\mathrm{C}^{-} = \mathrm{A} \cap \mathrm{D}^{-}\) are non-empty; these are chambers. Moreover, the hyperplane \(\mathbf{L}\) meets \(\mathrm{D}_{\mathfrak{N}}(\mathrm{F}) = \mathrm{D}_{\mathfrak{N}}(\mathrm{C}^{+})\) ; Prop. 4 shows that \(\mathbf{L}\) is a wall of \(\mathrm{C}^{+}\) and that
\end{enumerate}

F, which meets \(L \cap D_{\mathfrak{N}}(F)\) , is a face of \(C^{+}\) ; similarly, F is a face of \(C^{-}\) . Finally, let C be a chamber of which F is a face, and suppose for example that \(D^{+} = D_{L}(C)\) ; by Prop. 4, the set \(D_{\mathfrak{N}}(C)\) meets F, and hence is equal to \(D_{\mathfrak{N}}(F)\) , and we have

\[
C = D _ {\mathfrak {H}} (C) = D _ {L} (C) \cap D _ {\mathfrak {N}} (C) = D ^ {+} \cap D _ {\mathfrak {N}} (F) = C ^ {+}.
\]

